\subsection{Домашние задачи на четвёртую неделю}
\begin{problem}
    Найти математическое ожидание числа бросаний кости до первого выпадения двух шестёрок подряд.
\end{problem}
\begin{solution}
    Пусть $P(n)$~---~вероятность того, что последовательность из $n$ бросков не содержит двух шестёрок подряд.
    Тогда $P(0) = P(1) = 1$.
    При $n \geq 2$ справедливо следующее рекуррентное соотношение:
    \[
        P(n) = \dfrac{5}{6}P(n - 1) + \dfrac{5}{36}P(n - 2).
    \]
    Теперь пусть $p(n)$~---~вероятность того что игра закончится на $n$-ом броске.
    Тогда, поскольку на $n - 2$-ом броске не могла выпасть шестёрка, то для $p(n)$ справедливо следующее выражение:
    \[
        p(n) = P(n - 3)\cdot \dfrac{5}{6} \cdot \dfrac{1}{6^2} = \dfrac{5P(n - 3)}{216}, n \geq 3
    \]
    Математическим ожиданием числа бросков кости будет являться $\sum\limits_{n = 0}^{+\infty}np(n)$.
    Тогда
    \[
       \Ex = \sum\limits_{n = 0}^{+\infty} np(n) = \dfrac{2}{36} + \sum\limits_{n = 3} \dfrac{5nP(n - 3)}{216}.
    \]
    Найдём производящуюю функцию $G(z)$ для последовательности $P(n)$:
    \[
        G(z) = \sum\limits_{n = 0}^{+\infty} P(n)z^n = 1 + z + \sum\limits_{n = 2}^{+\infty} P(n)z^n = 1 + z + \dfrac{5z}{6}(G(z) - 1) + \dfrac{5z^2}{36}G(z).
    \]
    И тогда
    \[
        G(z)\biggr(1 - \dfrac{5}{6}z - \dfrac{5}{36}z^2\biggr) = 1 + \dfrac{z}{6}.
    \]
    То есть
    \[
        G(z) = \dfrac{1 + \dfrac{z}{6}}{1 - \dfrac{5}{6}z - \dfrac{5}{36}z^2} = \dfrac{36 + 6z}{36 - 30z - 5z^2}.
    \]
    Найдём $\sum nP(n)$:
    \[
        zG'(z) = \sum\limits_{n = 0}^{+\infty} nP(n)z^n.
    \]
    Теперь мы можем выразить $\Ex$ через производящую функцию и получить численный ответ:
    \[
        \Ex = \dfrac{1}{18} + \dfrac{5}{216}\sum nP(n) + \dfrac{15}{216}G(1) = \dfrac{12 + 5G'(1) + 15G(1)}{216}.
    \]
    Посчитаем $G(1):$
    \[
        G(1) = \dfrac{36 + 6z}{36 - 30z - 5z^2}\biggr|_{z = 1} = 42.
    \]
    Посчитаем $G'(1)$:
    \[
        G'(1) = \dfrac{6(-5z^2 - 30z + 36) - (6z + 36)(-10z - 30)}{(-5z^2 - 30z + 36)^2}\biggr|_{z = 1} = 6 + 42 \cdot 40.
    \]
    И получим
    \[
        \Ex = \dfrac{12 + 5 \cdot (6 + 42 \cdot 40) + 15 \cdot 42}{216} = \dfrac{(200 + 15 + 1)42}{216} = 42.
    \]
\end{solution}
\begin{problem}
    Найти математическое ожидание и дисперсию равномерно распределённой в $[a, b]$ случайной величины.
\end{problem}
\begin{solution}
    Плотность равномерно распределённой в $[a, b]$ случайной величины $\xi$ равна \[f_\xi = \dfrac{1}{b - a}I_{[a, b]}(x).\]
    Тогда
    \[
        \Ex \xi = \int_a^b \dfrac{x}{b - a}dx = \dfrac{(b^2 - a^2)}{2(b - a)} = \dfrac{a + b}{2}.
    \]
    И
    \[
        \D \xi = \int_a^b \dfrac{x^2}{b - a}dx - \dfrac{(a + b)^2}{4} = \dfrac{b^3 - a^3}{3(b - a)} - \dfrac{(a + b)^2}{4} = \dfrac{b^2 + ab + a^2}{3} - \dfrac{a^2 + 2ab + b^2}{4} = \dfrac{(a - b)^2}{12}
    \]
\end{solution}
\begin{problem}
    Известно, что случайная величина $\xi$ принимает значения из отрезка $[a, b]$. В каких пределах могут лежать $\Ex \xi$ и $\D \xi$? Укажите их как можно точнее.
\end{problem}
\begin{solution}
    Очевидно, что поскольку $a \leq \xi \leq b$, то $a = \Ex a \leq \Ex \xi \leq \Ex b = b$.
    Несложно заметить, что $\D (a\xi + b) = a^2\D\xi$ и поэтому мы можем без ограничения общности считать что случайная величина принимает значения в отрезке $[0, 1]$.
    Но тогда $\xi^2 \leq \xi$, а следовательно $\D \xi = \Ex \xi^2 - (\Ex \xi)^2 \leq \Ex \xi - (\Ex \xi)^2 \leq 1/4$.
    Нижняя оценка достигается на константе, верхняя -- на бернуллиевской случайной величине с параметром $1/2$.
    Используя ранее выписанное тождество получим, что для произвольного отрезка границами для дисперсии являются $0$ и $\dfrac{(b - a)^2}{4}$.
\end{solution}
\begin{problem}
    Случайная величина $\xi$ может принимать значения $-2, -1, 0, 1, 2$. Известно, что $\Ex \xi = \Ex \xi^3 = 0, \Ex \xi^2 = 1, \Ex \xi^4 = 2$.
    Найдите распределение $\xi$.
\end{problem}
\begin{solution}
    В силу дискретности случайных величин верна следующая система уравнений:
    \begin{equation*}
        \begin{cases}
            -2\mathbb{P}(\xi = -2) - \mathbb{P}(\xi = - 1) + \mathbb{P}(\xi = 1) + 2\mathbb{P}(\xi = 2) = 0 \\
            -8\mathbb{P}(\xi = -2) - \mathbb{P}(\xi = -1) + \mathbb{P}(\xi = 1) + 8\mathbb{P}(\xi = 2) = 0 \\
            4\mathbb{P}(\xi = -2) + \mathbb{P}(\xi = -1) + \mathbb{P}(\xi = 1) + 4\mathbb{P}(\xi = 2) = 1 \\
            16\mathbb{P}(\xi = -2)   + \mathbb{P}(\xi = -1) + \mathbb{P}(\xi = 1) + 16\mathbb{P}(\xi = 2) = 2
        \end{cases}
    \end{equation*}
    Из первых двух уравнений получим: $\mathbb{P}(\xi = 2) = \mathbb{P}(\xi = -2)$. Но, тогда, и $\mathbb{P}(\xi = -1) = \mathbb{P}(\xi = 1)$. Вычитая из 4 уравнения 3-е получим $24\mathbb{P}(\xi = 2) = 1$, то есть $\mathbb{P}(\xi = 2) = \mathbb{P}(\xi = -2) = \dfrac{1}{24}$. Значит $\mathbb{P}(\xi = -1) = \mathbb{P}(\xi = 1) = \dfrac{1}{3}$. В силу нормировки $\mathbb{P}(\xi = 0) = 1 - \dfrac{2}{3} - \dfrac{1}{12} = \dfrac{1}{4}$.
\end{solution}
\begin{problem}
    Пусть $\xi \sim N(0, 1)$.
    Вычислите $\Ex \xi^k$ для всех натуральных $k$.
\end{problem}
\begin{solution}
    Поскольку $\xi \sim N(0, 1)$, то $\xi$ имеет плотность $f_\xi(x) = \dfrac{1}{\sqrt{2\pi}}e^{-\frac{x^2}{2}}$. Пусть $k \in \N, k \geq 2$. Тогда
    \begin{multline*}
        \mathbb{E}\xi^k = \int_\mathbb{R}x^k\dfrac{1}{\sqrt{2\pi}}e^{-\frac{x^2}{2}}dx =  -\dfrac{1}{\sqrt{2\pi}}\int_\mathbb{R}x^{k - 1}de^{-\frac{x^2}{2}} = \\ = -\dfrac{1}{\sqrt{2\pi}}\biggr(x^{k - 1}e^{-\frac{x^2}{2}}\big|_{-\infty}^{+\infty} - (k - 1)\int_\mathbb{R}x^{k - 2}e^{-\frac{x^2}{2}}dx\biggr) = (k - 1)\mathbb{E}\xi^{k - 2}
    \end{multline*}
    В силу нормировки плотности, то есть \[
                                             \int_{-\infty}^{+\infty}f_\xi(x)dx = 1
    \]
    получается, что если $2 \mid k$, то $\mathbb{E}\xi^k = (k - 1)!!$. В случае $2 \nmid k$ мы получаем, что $\mathbb{E}\xi \propto \int_{-\infty}^{+\infty}de^{-\frac{x^2}{2}} = 0$ и $\mathbb{E}\xi^k = 0$.
\end{solution}
\begin{problem}
    Случайная величина $\xi$ имеет гамма-распределение $\Gamma(n, \lambda)$ с плотностью $f(x) = cx^{n - 1}e^{-\lambda x}I_{\R_+}(x)$.
    Найти $c, \Ex \xi$ и $\D \xi$.
\end{problem}
\begin{solution}
    Будем считать, что $\lambda > 0$, поскольку иначе интеграл расходится и $f$ не может быть плотностью распределения. Поскольку плотность должна быть нормирована, то \[
                                                                                                                                                                            \int_{-\infty}^{+\infty}fdx = 1
    \]
    То есть \[
                \int_0^{+\infty}c(n, \lambda)x^{n - 1}e^{-\lambda x}dx = 1
    \]
    Поскольку
    \begin{multline*}
        \dfrac{1}{c(n, \lambda)} = \int_{\mathbb{R}_+} x^{n - 1}e^{-\lambda x}dx = \\ =  -\dfrac{1}{\lambda}\biggr(x^{n - 1}e^{-\lambda x}\big|^{+\infty}_0 - (n - 1)\int_{\mathbb{R}_+}x^{n - 2}e^{-\lambda x}dx\biggr) = \\ = \dfrac{n - 1}{\lambda c(n - 1, \lambda)}
    \end{multline*}
    То есть $c(n, \lambda) = \dfrac{\lambda}{n - 1}c(n - 1, \lambda)$. Значит $c(n, \lambda) = \dfrac{\lambda^{n - 1}}{(n - 1)!}c(1, \lambda)$. А вот это уже можно посчитать: \[
                                                                                                                                                                                   \int_{\mathbb{R}_+} e^{-\lambda x}dx = -\dfrac{1}{\lambda}e^{-\lambda x}\big|_0^{+\infty} = \dfrac{1}{\lambda}
    \]
    Значит $c(n, \lambda) = \dfrac{\lambda^n}{(n - 1)!} = \dfrac{\lambda^n}{\Gamma(n)}$. \\ Теперь посчитаем $\mathbb{E}\xi$:
    \begin{equation*}
        \mathbb{E}\xi = \int_{\mathbb{R}} xf(x)dx = c(n, \lambda) \int_{\mathbb{R}_+} x^{(n + 1) - 1}e^{-\lambda x}dx = \dfrac{c(n, \lambda)}{c(n + 1, \lambda)} = \dfrac{\Gamma(n + 1)}{\Gamma(n)}\dfrac{\lambda^n}{\lambda^{n + 1}} = \dfrac{n}{\lambda}
    \end{equation*}
    И $\mathbb{D}\xi$:
    \begin{equation*}
        \mathbb{D}\xi = \mathbb{E}\xi^2 - \big(\mathbb{E}\xi\big)^2 = \dfrac{c(n, \lambda)}{c(n + 2, \lambda)} - \dfrac{n^2}{\lambda^2} = \dfrac{(n + 1)n - n^2}{\lambda^2} = \dfrac{n}{\lambda^2}
    \end{equation*}
\end{solution}
\begin{problem}
    Пусть $Z$~---~расстояние между двумя случайно и независимо выбранными точками в круге радиуса $r$.
    Вычислить $\Ex Z^2$.
\end{problem}
\begin{solution}
    Пусть выбраны точки $P = (x_1, y_1)$ и $Q = (x_2, y_2)$. Тогда $Z^2(P, Q) = (x_1 - x_2)^2 + (y_1 - y_2)^2$. Пространством является $(B^2\times B^2, \mathcal{B}, \mathcal{L}^4)$. Тогда,
    \begin{multline*}
        \mathbb{E}Z^2\cdot \mathcal{L}^4(B^2 \times B^2) = \iiiint_\Omega (x_1 - x_2)^2 + (y_1 - y_2)^2dx_1dx_2dy_1dy_2 = \iiiint_\Omega x_1^2 + x_2^2 + y_1^2 +y_2^2dx_1dx_2dy_1dy_2 = \\ =  \int_0^{2\pi}d\varphi\int_0^{2\pi}d\psi\int_0^rda\int_0^r (a^2 + b^2)abdb = 4\pi^2 \int_0^r \dfrac{a^3r^2}{2} + \dfrac{ar^4}{4}da = 4\pi^2\cdot \dfrac{r^6}{4} = \pi^2r^6
    \end{multline*}
    После нормировки на $\pi^2 r^4$ получаем искомое матожидание $\mathbb{E}Z^2 = r^2$.
\end{solution}
\begin{problem}
    Пусть $\xi$ имеет конечную дисперсию.
    Найдите какие константы $a$ и $b$ минимизируют $\phi(a) = \Xi (\xi - a)^2$ и $\psi(b) = \Ex |\xi - b|$.
\end{problem}
\begin{solution}
    \begin{itemize} \leavevmode
    \item Раскроем скобочки: \[\varphi(a) = \int_\Omega(\xi - a)^2d\mathbb{P} = \int_\Omega\xi^2d\mathbb{P} - 2a\int_\Omega\xi d\mathbb{\mathbb{P}} + a^2  = \mathbb{E}\xi^2 - 2a\mathbb{E}\xi + a^2\]
    Очевидно, что $\varphi(a)$ минимальна в $a = -\dfrac{-2\mathbb{E}\xi}{2} = \mathbb{E}\xi$. В таком случае $\varphi(\mathbb{E}\xi) = \mathbb{D}\xi$.
    \item Пусть $t \in \mathbb{R}$~---~такое число, что $\mathbb{P}(\xi \leq t) \geq 0.5$ и $\mathbb{P}(\xi \geq t) \geq 0.5$. Тогда $t$~---~искомое число, минимизирующее $\psi(b)$. Пусть есть некоторое $t^* > t$. Тогда \begin{multline*}
                                                                                                                                                                                                                                \mathbb{E}(|\xi - t^*| - |\xi - t|) = \int_\Omega |\xi - t^*| - |\xi - t|d\mathbb{P} = \\ = \int_{\mathbb{P}(\xi > t^*)} t - t^*d\mathbb{P} + \int_{\mathbb{P}(t^* \geq \xi > t )}t^* + t - 2\xi d\mathbb{P} + \int_{\mathbb{P}(t \geq \xi)} t^* - td\mathbb{P} \geq \\ \geq (t - t^*)\mathbb{P}(\xi > t) + (t^* - t)\mathbb{P}(\xi \leq t)  = (t^* - t)(2\mathbb{P}(\xi \leq t) - 1) \geq 0
    \end{multline*}
    Если же $t^* < t$, то мы можем поменять $t$ и $t^*$ в предыдущем выражении местами, получить схожую оценку и в силу монотонности меры получить искомое.
    \end{itemize}
\end{solution}